% =====================================================================
% CHAPTER 4: System Development and Implementation (new 9-section
% department template, thesis-or-project-template_1745300785.pdf).
% Assembled per thesis/RESTRUCTURE_PLAN_V2.md from old
% mainmatter/ch7.tex sections 7.3 (Selection of Parameters, Apparatus
% and Components), 7.4 (Optimization and Configuration Selection), and
% 7.7 (Principle of Operation), with fresh 4.1 Introduction and
% Summary prose. See the % SOURCES / DROPPED / MOVED-AWAY block at the
% end of this file for the full accounting.
% =====================================================================
\chapter{System Development and Implementation}
\label{ch:implementation}

\section{Introduction}
A design on paper only becomes useful once it is written as code,
trained until it stops improving, and made fast enough to run on the
hardware it was meant for. This chapter is about that step. It turns
StereoLite from a fixed architecture into a working system and records,
in one place, everything needed to rebuild it: the software that trains
and evaluates the network, the exact settings the final model uses, the
work done to bring one forward pass inside the target speed and memory
budget, the computing and camera hardware involved, and what actually
happens inside the network when a rectified image pair goes in and a
disparity map comes out.

The emphasis here is on what was built, not on the numbers it produced;
those belong with the results. Each value that matters, how the network
is sized, how it is trained, and what is changed to run it in eight-bit
on an edge board, is pinned down once so that every result reported
afterwards rests on a configuration a reader can see in full and
reproduce.

\FloatBarrier
\section{Software Development}
\label{sec:software_development}
The network is implemented in PyTorch 2.11 with CUDA 12.8 on Python
3.10. We organized the codebase around scripted drivers rather than
notebooks: a separate script trains the production model, separate
harnesses run the encoder, refinement, and loss ablations under a
matched protocol, and separate drivers run the SceneFlow held-out
evaluation and the zero-shot Middlebury 2014 cross-domain evaluation.
Each driver writes its full configuration and its metrics to a metadata
file next to its outputs, so a reported number can be traced back to
the exact run that produced it.

Two cloud platforms carry the compute load. Full-data training and the
production run described in Section \ref{sec:network_implementation}
ran on Modal-orchestrated cloud A100 jobs, while ablations and
cross-domain evaluations ran on cloud T4 instances, which are cheap
enough to run after every design change without straining the project
budget. An earlier development pass on the edge-tier chassis used a
self-contained Kaggle notebook with data-parallel training across two
T4 GPUs; the production run reported in this thesis moved to the
single-A100 Modal path once the architecture and loss were fixed, since
a single larger GPU trains the full SceneFlow corpus faster than two
smaller ones split across processes. Keeping training, ablation, and
evaluation as separate scripted drivers, each logging its own metadata,
is what let every reported number trace back to a specific run rather
than to a notebook cell that may since have changed.

\FloatBarrier
\section{Network Implementation}
\label{sec:network_implementation}
The design chapter argued for each architectural family; this section
records the exact values the production model uses. The backbone is a
YOLO26s encoder truncated after its first seven blocks. The recurrent
hidden state carries 32 channels, and the refinement head that reads it
carries 48. The tile initialization module builds its correlation
volume from eight groups over 24 disparity hypotheses at one sixteenth
resolution, and the per-iteration local correlation looks up a
half-range of two pixels around the current disparity estimate. The
refinement schedule runs two iterations at one sixteenth resolution,
three at one eighth, and three at one quarter, and the geometry
encoding volume is narrowed to a band of plus or minus 16 disparities
around the tile estimate rather than spanning the full disparity range.
Table \ref{tab:arch_config} of Appendix~\ref{app:specs} gives the complete
configuration record of the trained checkpoint, verbatim, alongside
these headline values.

Table \ref{tab:train_config} lists the training-side configuration that
Chapter \ref{ch:design} selected and justified: the optimizer, schedule,
batch size, input protocol, augmentation set, training data, validation
protocol, and training hardware. Every value in the table is the result
of a controlled ablation, and the design-evidence audit of
Chapter \ref{ch:design} traces each choice to its justification.
Section \ref{sec:training_objective} below fixes the
training objective and the protocol narrative behind this table in
full.

\begin{table}[H]
\centering
\caption{Training configuration of the production run (run 20260704).}
\label{tab:train_config}
\small
\begin{tabular}{ll}
\toprule
Element & Value\\
\midrule
Optimizer & AdamW, weight decay $10^{-5}$\\
Schedule & OneCycle, 60{,}000 steps, peak LR $8\times 10^{-4}$\\
Batch size / precision & 32 / bfloat16 mixed\\
Input protocol & random native $384\times 640$ co-located crops\\
Augmentation & asymmetric color jitter, right-only erase, random scale\\
Training data & full SceneFlow finalpass, 35{,}454 pairs\\
Validation & fixed 400-pair FlyingThings3D subset, every 1k steps\\
Hardware & one NVIDIA A100 80 GB (cloud)\\
\bottomrule
\end{tabular}
\end{table}

\FloatBarrier
\section{Optimization and Configuration Selection}
\label{sec:optimization}
Reaching the latency target required optimizing the network's
inference graph, not just training it. We applied six changes: we
rewrote the cost-volume construction loops to use padded zero-copy
views instead of allocating and zeroing a feature map for every
disparity bin; we reduced the grid-sample calls from 48 to eight by
batching the lookups and caching the sampling grids; we hoisted the
static channels of the recurrent input, which are re-convolved every
iteration, to run once per scale, following the precomputed-context
trick of RAFT-Stereo \cite{lipson2021raftstereo}; we fused the gate
convolutions and the four output heads, which share inputs; we removed
dead inference-path code; and we narrowed the geometry volume to a band
around the tile disparity. We proved each change equivalent before
adopting it, either bitwise or through a matched A/B comparison, and
together they cut the full-precision latency by a factor of 1.74.

Figure \ref{fig:export_pipeline} lays the deployment path out: the four
graph optimizations act on the full-precision model before export,
while the three operator swaps prepare the graph for eight-bit
calibration.

\begin{figure}[!htbp]
\centering
\includegraphics[width=\textwidth]{fig_4_5_export_preview}
\caption{Export pipeline from trained checkpoint to INT8 TensorRT engine.}
\label{fig:export_pipeline}
\end{figure}

The same pass also had to clear three operators that block eight-bit
deployment: the unfold operation in convex upsampling, group
normalization on some accelerator toolchains, and three-dimensional
convolution on neural processing units. We swapped each for a
quantization-friendly equivalent, then exported and calibrated the
resulting graph to an eight-bit TensorRT engine for the Jetson Orin
Nano. Table \ref{tab:efficiency_findings} summarizes the six changes and
their equivalence status.

\begin{table}[!htbp]
\centering
\caption{Inference optimizations and equivalence status measured at $384\times 640$ on RTX 3050.}
\label{tab:efficiency_findings}
\small
\begin{tabular}{p{5.6cm}p{4.2cm}l}
\toprule
Optimization & Saving & Equivalence\\
\midrule
Zero-copy padded views in cost-volume loops & 5 to 15\% latency & bitwise\\
Batched correlation lookups, cached grids & 3 to 8\% latency & identical math\\
Static GRU input channels hoisted per scale & 8 to 10\% of MACs & exact\\
Fused gate convolutions and output heads & 1 to 3\% latency & bitwise\\
Dead inference-path code removed & memory + launches & exact\\
Narrow GEV band ($\pm 16$ around tile $d$) & 25 to 30\% of MACs & matched A/B\\
\midrule
Combined & fp32 61.4 ms vs 106.7 ms base ($1.74\times$) & \\
\bottomrule
\end{tabular}
\end{table}

\FloatBarrier
\section{Training Objective and Protocol}
\label{sec:training_objective}
The optimizations above act on a network that already has to be trained
correctly in the first place. This section fixes what the network
minimizes during training and the data protocol under which that
objective is applied, the two pieces of the design methodology that sit
upstream of the parameter values already listed in Table
\ref{tab:train_config}.

\subsection{Training Objective}
The network trains under the multi-scale objective of \eqref{eq:loss}. A
pixel counts as valid when its ground-truth disparity is finite and lies
between 0 and 320 pixels. Each lower-scale prediction is resized to full
resolution and multiplied by its scale factor before comparison. The
objective weights an L1 term at every scale, heaviest at full
resolution, adds a matching L1 term on the geometry-volume output, and
augments the full-resolution prediction with a Sobel
gradient-consistency term, a threshold hinge stack, a KITTI-style
D1-relative hinge \cite{menze2015kitti}, an edge-aware smoothness term, and a gated
slant-supervision term that trains the plane slopes only where they are
trustworthy. The slant term pairs with plane rendering: it makes the
slopes accurate enough for the rendering step to sharpen boundaries.

We chose this composition by ablation rather than guesswork. The loss
study tested nine formulations: plain L1
under-serves the outlier metrics, hinge terms on a single threshold
overfit that threshold, and the threshold stack combined with the D1
hinge was the only formulation to win both the end-point error and the
outlier rates. The production objective descends from that winner,
keeping the gradient term for its sub-pixel benefit and holding the
smoothness term to a small weight so it regularizes without blurring.

\begin{equation}
\begin{aligned}
\mathcal{L} \;=\;& 1.00\,\mathcal{L}_{1}(d_{\mathrm{full}})
 + 0.50\,\mathcal{L}_{1}(d_{1/2})
 + 0.30\,\mathcal{L}_{1}(d_{1/4})
 + 0.20\,\mathcal{L}_{1}(d_{1/8})
 + 0.10\,\mathcal{L}_{1}(d_{1/16})\\
&+ 0.15\,\mathcal{L}_{1}(d_{\mathrm{gev}})
 + 0.50\,\mathcal{L}_{\mathrm{grad}}
 + 0.20\,\mathcal{L}_{\mathrm{thr}}
 + 0.20\,\mathcal{L}_{\mathrm{D1}}
 + 0.02\,\mathcal{L}_{\mathrm{smooth}}
 + 0.30\,\mathcal{L}_{\mathrm{slant}}
\end{aligned}
\label{eq:loss}
\end{equation}
where $\mathcal{L}_1$ is the masked mean absolute disparity error,
$\mathcal{L}_{\mathrm{grad}}$ a Sobel gradient-consistency term,
$\mathcal{L}_{\mathrm{thr}}$ a hinge over the error thresholds
$\{0.5, 1, 2, 3\}$ pixels, $\mathcal{L}_{\mathrm{D1}}$ the D1-relative
hinge, $\mathcal{L}_{\mathrm{smooth}}$ an edge-aware smoothness term, and
$\mathcal{L}_{\mathrm{slant}}$ the gated slant supervision.

\begin{figure}[!htbp]
\centering
\includegraphics[width=\textwidth]{fig_3_8_supervision_preview}
\caption{Six disparity scales supervised with L1, gradient consistency, and threshold terms.}
\label{fig:supervision}
\end{figure}

\FloatBarrier
\subsection{Datasets and Training Protocol}
This subsection fixes the protocol applied on top of the SceneFlow
finalpass corpus, the thesis's training and in-domain test set, and the
project's own rig captures used only for qualification. Optimization uses
AdamW with a OneCycle schedule over 60,000 steps at a peak learning rate
of $8 \times 10^{-4}$, a batch of 32, and bfloat16 mixed precision, on a
single A100 GPU. Input is normalized to the unit interval at the model
boundary. Table \ref{tab:train_config} lists this configuration in full;
the paragraphs below detail the two parts a reader is most likely to want
to reproduce, the data split and the crop-and-augment pipeline.

\paragraph{Data split.}
Table \ref{tab:data_split} gives the exact split. SceneFlow finalpass is
the union of three rendered subsets, FlyingThings3D, Monkaa, and Driving,
whose training halves sum to the 35,454 pairs the production model trains
on. Only FlyingThings3D carries an official held-out test split, 4,370
pairs, which serves as the in-domain benchmark and is never seen during
training. In-training validation is not a fourth pool of data: it is a
fixed 400-pair subset drawn from that test split, evaluated every
thousand steps at the native $960\times 540$ axis so that progress is
tracked against the same distribution the final number is reported on.
The project's own stereo rig contributes 997 real indoor pairs, and its
split is deliberately degenerate: none of them train or fine-tune the
network at any point, so all 997 sit on the test side as a zero-shot
real-camera qualification set, of which three scenes are carried through
to the metric reconstruction of Chapter~\ref{ch:results}.

\begin{table}[!htbp]
\centering
\caption{Dataset split used for training, validation, and evaluation.}
\label{tab:data_split}
\small
\begin{tabular}{llrrr}
\toprule
Source & Subset & Train & Val & Test\\
\midrule
SceneFlow finalpass & FlyingThings3D & 22{,}390 & 400$^{\dagger}$ & 4{,}370\\
SceneFlow finalpass & Monkaa & 8{,}664 & & \\
SceneFlow finalpass & Driving & 4{,}400 & & \\
\midrule
SceneFlow finalpass & total & 35{,}454 & 400$^{\dagger}$ & 4{,}370\\
\midrule
Stereo rig & indoor captures & 0 & 0 & 997\\
\bottomrule
\end{tabular}

\smallskip
{\footnotesize $^{\dagger}$The 400-pair validation set is a fixed subset
drawn from the FlyingThings3D test split, evaluated every 1{,}000 steps;
it is not additional data.}
\end{table}

\paragraph{Cropping.}
Training samples a random $384\times 640$ window, co-located at the same
coordinates in both the left and right view, from each native $960\times
540$ frame. Cropping at the native resolution rather than downscaling the
whole frame first is a deliberate choice: a disparity is a distance in
pixels, so shrinking the image also shrinks every disparity, and a
network trained on downscaled frames learns to predict a compressed
disparity range that no longer matches full-resolution input. The native
crop keeps disparity magnitudes at their true scale while still fitting
the memory budget, and the ablation of Chapter~\ref{ch:design} found this
native-crop protocol substantially ahead of the legacy
downscale-then-crop protocol; Figure~\ref{fig:input_protocol} contrasts
the two. Cropping the same window from both views is what preserves the
epipolar correspondence: shifting the crops relative to each other would
silently offset every disparity by the shift, so the window is always
identical in the two images.

\paragraph{Augmentation.}
Three augmentations, drawn from the OpenStereo recipe, are applied to the
training crop only; Figure~\ref{fig:augmentation} shows each acting on a
real pair, and each answers a specific failure mode of a
synthetic-only model on real cameras. \emph{Asymmetric color jitter}
perturbs brightness, contrast, and saturation each within
$[0.6, 1.4]$, and one time in five it draws different parameters for the
left and right image rather than the same ones. A real low-cost stereo
rig has two separate sensors that never expose or white-balance
identically, so the two views differ photometrically; jittering them
asymmetrically forces the matcher to rely on structure rather than on an
assumption that corresponding pixels share an exact color, which is what
lets the SceneFlow-trained model survive on the rig. \emph{Right-image
random erasing} blanks one or two rectangles of 50 to 100 pixels in the
right view only, filling them with the image mean. Because the erased
content has no match in the left view, this manufactures occlusion on
demand and trains the network to infer disparity from context where a
direct match is unavailable, exactly the situation at object boundaries
and image borders. \emph{Random scaling with anisotropic stretch}
resizes the crop by a factor of $2^{U(-0.2,0.4)}$, with an extra
horizontal stretch of $2^{U(-0.2,0.2)}$, and rescales the disparity by
the horizontal factor before recropping to $384\times 640$. Scaling
widens the range of object sizes and, crucially, of disparity magnitudes
the network sees, so it does not overfit to the specific scale statistics
of the synthetic set; the horizontal stretch changes disparity directly,
which is why the disparity target is multiplied by the same horizontal
factor. Evaluation applies none of these, so every reported number
measures the network on unperturbed data.

\begin{figure}[!htbp]
\centering
\includegraphics[width=\textwidth]{fig_4_3_augmentation}
\caption[Training augmentations applied to a native co-located crop.]{Training augmentations applied to a native $384\times 640$ co-located crop of a real rectified pair (Middlebury 2014 Motorcycle \cite{scharstein2014middlebury}): asymmetric color jitter, right-only random erasing (boxes outlined), and random scale with horizontal stretch. The eraser and the two views' color jitter act independently on the right image; scaling rescales the disparity target by the horizontal factor.}
\label{fig:augmentation}
\end{figure}

\begin{figure}[!htbp]
\centering
\includegraphics[width=\textwidth]{fig_3_10_input_protocol}
\caption{Native-crop sampling (left) vs downscale-then-crop (right) training protocols.}
\label{fig:input_protocol}
\end{figure}

\FloatBarrier
\section{Hardware and Embedded Deployment}
\label{sec:hw_deployment}
The apparatus divides into compute and sensing. On the compute side,
training runs on a single cloud A100 GPU with 80 gigabytes of memory.
Ablations and evaluations run on a cloud T4, chosen for its low cost,
which keeps a controlled study practical after every design change. The
local latency benchmark runs on an RTX 3050 laptop GPU with roughly 4
gigabytes of memory, representing the low end of the deployment memory
budget. The intended deployment target is the Jetson Orin Nano class of
edge accelerator, on which we measured the optimized eight-bit model
directly.

On the sensing side, real-camera pairs come from a calibrated low-cost
binocular rig that streams a 2560 by 720 side-by-side image, split into
a rectified 1280 by 720 pair per frame. This rig supplies the real-data
qualification of the trained model: the zero-shot qualitative pairs, the
real-camera agreement protocol, and the three-dimensional
reconstructions it supports. Housing both
the training-side compute and the sensing-side rig in the same section
keeps the apparatus that produced every downstream number in one place,
ahead of the experimental protocols built on top of it.

\FloatBarrier
\section{Principle of Operation}
\label{sec:principle_of_operation}
In deployment the system runs as a fixed feed-forward pipeline. A
rectified left and right pair enters at the inference resolution of 384
by 640 pixels, normalized to the unit interval, and passes through the
blocks fixed in Chapter \ref{ch:design}. The shared encoder processes
both views in a single batched pass and emits matching features at one
quarter, one eighth, and one sixteenth resolution, while a separate
context encoder reads the left image alone and produces the context
maps that feed the recurrent state.

Matching begins at the coarsest scale. The tile initialization module
builds the group-wise correlation volume of \eqref{eq:gwc} at one
sixteenth resolution, regularizes it with a small three-dimensional
network, and reads off the initial disparity and confidence through the
soft-argmax of \eqref{eq:softargmax_init}; the plane slopes start at
zero, and the hidden state starts from the coarsest context map. The
recurrent refinement then runs its fixed schedule: two iterations at
one sixteenth, three at one eighth, and three at one quarter resolution.
At each iteration the right features are warped by the current
disparity, a local correlation is evaluated over five offsets, and the
convolutional gated recurrent unit of \eqref{eq:convgru} updates the
hidden state; gated residuals adjust the disparity, the plane slopes,
and the confidence. Between scales, the plane-aware propagation of
\eqref{eq:plane_prop} carries the tile state to the finer grid along its
local slanted-surface hypothesis rather than by plain interpolation.

At one quarter resolution, the narrow-band geometry encoding volume
produces its independent proposal per \eqref{eq:gev_expect}, and the
fail-soft gate of \eqref{eq:failsoft} blends that proposal into the
recurrent estimate only where the learned weight favors it. This
corrects regions the tile state resolved poorly, without disturbing the
ones it resolved well. The refined quarter-resolution state is rendered
from its tile planes with an edge gate and upsampled to full resolution
by two learned convex-combination stages, each output subpixel a
softmax-weighted mixture of its three-by-three coarse neighborhood. The
result is a dense full-resolution disparity map.

Disparity becomes depth through the calibrated relations of
\eqref{eq:backprojection}: depth follows $Z = fB/d$, and each valid
pixel optionally back-projects to a metric three-dimensional point in
the camera frame, colored by its left-image pixel. On the project's
rig, the calibrated focal length and baseline scale to the inference
resolution by the width ratio, so the point cloud is metric without
further tuning. This single forward pass, from image pair to disparity
map and optional colored point cloud, is the operation this thesis
measures throughout.

\FloatBarrier
\section{Summary}
This chapter traced the path from the network architecture to a
trained, deployable system. The codebase separates training, ablation,
and evaluation into scripted drivers that each log their own
configuration and metrics, running on Modal-orchestrated cloud A100 and
T4 GPUs. The production model fixes a truncated YOLO26s backbone, a
32-channel recurrent state, an eight-group correlation volume, and a
two, three, three iteration schedule, and this configuration is
recorded in full. A six-change, equivalence-proven optimization pass
cut inference latency by a factor of 1.74 and cleared the operators
that block eight-bit deployment, producing a calibrated TensorRT engine
for the Jetson Orin Nano. Training itself runs under an eleven-term
multi-scale objective built around a threshold-stack and D1-relative
hinge, applied to random native-resolution crops of the SceneFlow
corpus with a three-part augmentation recipe and a sixty-thousand-step
OneCycle schedule. The compute and sensing apparatus spans a cloud A100
for training, a cloud T4 for ablations, an RTX 3050 for local latency
benchmarking, the Jetson Orin Nano deployment target, and a calibrated
low-cost stereo rig for real-camera qualification. The principle of
operation section walked through the resulting forward pass, from a
rectified image pair to a full-resolution disparity map and an optional
metric point cloud.

% SOURCES: mainmatter/ch7.tex sections 7.2 (Training Objective and
% Datasets and Training Protocol subsections -> new section "Training
% Objective and Protocol", added 2026-07-17 patch: eq:loss, fig:supervision,
% fig:input_protocol carried intact, numbers frozen), 7.3 (Selection of
% Parameters, Apparatus and Components -> split across Network
% Implementation and Hardware and Embedded Deployment), 7.4 (Optimization
% and Configuration Selection, verbatim numbers and table), and 7.7
% (Principle of Operation, unchanged aside from chapter cross-references).
% New prose: Introduction, Software Development (expanded from the compute
% half of the old 7.3 apparatus paragraph plus documented Modal/Kaggle
% infrastructure facts, no new numbers), the connective prose in Training
% Objective and Protocol, and Summary. Labels carried unchanged:
% tab:train_config, tab:efficiency_findings, eq:loss, fig:supervision,
% fig:input_protocol.
%
% PATCH NOTE (2026-07-17): during the parallel restructure, ch3 and ch4
% each assumed the other owned old 7.2's training-objective block, so it
% was orphaned and \ref{eq:loss} dangled book-wide. This patch resolves
% that by giving ch4 ownership, matching ch3's own MOVED-AWAY comment
% ("old ch7 7.2 subsection Training Objective -> implementation chapter
% (ch4)") and ch5's forward reference ("Chapter~\ref{ch:implementation}
% describes the crop, augmentation, and optimization protocol"). Checked
% ch3.tex and ch5.tex for duplication first: neither defines eq:loss,
% fig:supervision, or fig:input_protocol, and neither restates the crop/
% augmentation/optimizer prose in full (ch5 only names the corpus and
% explicitly defers the protocol here; ch3 only cross-references the
% ablation results tables by role). The new section therefore introduces
% no duplicate prose; it references Chapter 5 for the corpus facts
% (pair counts, resolution) instead of restating them, and reuses the
% already-present Table \ref{tab:train_config} rather than redefining it.
%
% Equation labels referenced but NOT owned here (assumed to live in
% Chapter 3 per the ownership map, since they originate in old
% mainmatter/ch3.tex "Design of StereoLite"): eq:gwc,
% eq:softargmax_init, eq:convgru, eq:plane_prop, eq:gev_expect,
% eq:failsoft, eq:backprojection. Figure/label referenced but NOT owned
% here (belongs to Chapter 5 per old section 7.6): fig:jetson,
% sec:testing_protocols. Cross-chapter labels assumed by convention
% (confirm at integration): ch:design (Chapter 3), ch:experimental
% (Chapter 5), ch:results (Chapter 6).
% DROPPED: nothing substantive. The old 7.3 sentence pointing forward
% to "the real-camera agreement protocol of Section
% \ref{sec:testing_protocols}" and "the three-dimensional
% reconstructions of Chapter 8" is kept as a forward pointer in Section
% \ref{sec:hw_deployment}, renumbered to the new chapter labels.
% MOVED-AWAY: old 7.5 "Review and Feedback" belongs to Chapter 3 (design
% assessment). Old 7.6 "Simulation and Experimental Setup" (fig:setup,
% fig:jetson) and 7.8 "Testing and Measurement" (sec:testing_protocols)
% belong to Chapter 5. Old 7.1/7.9 framing sentences were dissolved; any
% implementation-flavored fragments were folded into the new Introduction
% and Summary rather than kept as standalone sections.
